\documentclass[12pt]{article}
\usepackage[margin=0.75in]{geometry}
\usepackage{setspace}
\usepackage{lmodern}
\renewcommand{\labelitemi}{\ensuremath{\bullet}}
\RequirePackage[colorlinks,citecolor=blue,linkcolor=blue,urlcolor=blue]{hyperref}
\setstretch{1.08}

\begin{document}

\begin{center}
    \singlespacing
    {\Large\bf Identifying Rational Types in Unknown Environments under an Indirect Dynamic Mechanism}
    
    \vspace{1.2cm}
    
    {\large Humberto Bernal\footnote{Economist, professor of the economics program at the Universidad Colegio Mayor de Cundinamarca and Ph.D. in economics from the Universidad de los Andes, Colombia. Corresponding author. E-mail: \href{mailto:hbernal@uniandes.edu.co}{hbernal@uniandes.edu.co}, \href{mailto:hbernalc@universidadmayor.edu.co}{hbernalc@universidadmayor.edu.co}. ORCID: \href{https://orcid.org/0000-0002-4864-1726}{0000-0002-4864-1726}. All research is the sole responsibility of the author, received no external funding, and was developed as an independent academic initiative.}}
    
    \vspace{0.4cm}
    
    {\small Economics Program, Universidad Colegio Mayor de Cundinamarca; Ph.D. in Economics, Universidad de los Andes, Bogot\'a, Colombia.}
    
    \vspace{1.2cm}
    
    \textbf{Abstract}
\end{center}

\begin{abstract}
This paper develops a dynamic indirect mechanism to identify rational types in unknown environments characterized by observational pooling. The central theoretical result, the Rational Type Identification Theorem, proves that the principal's Bayesian beliefs converge almost surely to the agent's true latent type along prolonged interactions without direct revelation reports. This asymptotic identification is operationalized by the Mano de Dios-Guadalupe process, an endogenous maximum-entropy noise kernel that ensures full support across actions within a competing-risks architecture. A complementary Implementability Theorem guarantees execution in Perfect Bayesian Equilibrium. Microeconometric calibration and deep-learning simulations of Colombian kidnapping validate rapid type separation and policy feasibility.
\end{abstract}

\vspace{0.8cm}

\noindent \textbf{Keywords:} indirect mechanism design, dynamic screening, endogenous implementation noise, Bayesian learning, type identification \\
\noindent \textbf{JEL Classification:} C73, D82, K42
\vspace{0.8cm}

\noindent \textbf{Statements and Declarations:}
\begin{itemize}\small
    \item \textbf{Funding:} The author declares that no funds, grants, or other financial support were received for the preparation of this manuscript.
    \item \textbf{Conflict of Interest / Competing Interests:} The author declares that he has no conflict of interest. The author has no relevant financial or non-financial competing interests to disclose.
    \item \textbf{Data and Code Availability Statement:} The datasets generated and analysed during the current study, as well as the replication scripts and simulation code, are openly available in the public repository at \url{https://github.com/Donzhumber/CMP.git}.
    \item \textbf{Author Contributions:} Humberto Bernal is the sole author and contributed to all aspects of the conceptualization, mathematical proofs, econometric estimation, deep learning architecture, dynamic simulations, and writing of the manuscript.
    \item \textbf{Declaration on the Use of Generative AI:} In the preparation of this manuscript, the author used generative AI tools (\emph{Cursor} and \emph{Antigravity}) for language copy-editing and assistance with software syntax in Lean 4, Stata, and Python. The author validated all results and takes full responsibility for the paper's contents.
\end{itemize}

\end{document}
